% Guia do professor — Capítulo 18: Camada Limite Atmosférica
\section{Capítulo 18 --- Camada Limite Atmosférica}

\subsection*{Visão geral e ritmo}

O capítulo em que a turma descobre que mora dentro do objeto de
estudo. Três argumentos dimensionais (perfil log, orçamento de TKE,
$-5/3$) fazem a ponte perfeita com a formação em física; o ciclo
diurno costura tudo. \textbf{Ritmo: 2 aulas de 2\,h.}

\begin{itemize}
  \item \textbf{Aula 1} — o ``dia da CLA'' (desenho-mestre),
        encroachment e $z_i$ (Caso~1), camada superficial e o log
        (Caso~2, com energia eólica).
  \item \textbf{Aula 2} — TKE e seus termos, Kolmogorov com o
        espectro sintético (Caso~3), a noite estável e o jato de
        Blackadar (Caso~4).
\end{itemize}

\subsection*{Objetivos de aprendizagem}

\begin{enumerate}
  \item narrar o ciclo diurno da CLA (mistura, residual, estável) e
        seus perfis de $\theta_v$;
  \item prever $z_i(t)$ pelo encroachment a partir de $F_H$ e
        $\gamma$;
  \item ajustar o perfil log ($u_*$, $z_0$) a dados de torre e
        aplicá-lo a energia eólica;
  \item interpretar o orçamento de TKE e o espectro $-5/3$;
  \item explicar o jato noturno como oscilação inercial de Blackadar.
\end{enumerate}

\subsection*{Pré-requisitos a reativar}

Balanço de superfície e Bowen (Cap.~3); perfis de $\theta_v$ e
estabilidade não local (Cap.~5); círculos inerciais e Ekman
(Cap.~10); resfriamento radiativo noturno (Caps.~2, 6); $Ri$
(Cap.~5).

\subsection*{Armadilhas conceituais frequentes}

\begin{itemize}
  \item \textbf{``A CLA é fina, logo irrelevante''}: é onde vive TODO
        o vapor dos Caps.~4--7, o combustível do Cap.~14 e a poluição
        do Cap.~19 — e onde o vento vira recurso.
  \item \textbf{Perfil log como lei universal}: vale na camada
        superficial neutra; estabilidade curva o perfil
        (Monin--Obukhov). Não extrapolar o log até $z_i$.
  \item \textbf{Turbulência como ruído}: tem orçamento, cascata e
        espectro — é física estatística organizada, não sujeira.
  \item \textbf{Jato noturno por ``canalização''}: é inércia pura
        (Blackadar); o hodógrafo circular do Caso~4 desfaz o mito.
  \item \textbf{Camada residual esquecida}: o ar da tarde de ontem
        sobrevive à noite em cima — e é o que a térmica de amanhã
        reencontra (e o que guarda poluição envelhecida, Cap.~19).
\end{itemize}

\subsection*{Demonstração de baixo custo}

Timelapse de cumulus de bom tempo (o teto plano dos cumulus $=$
$z_i$); chaleira fumegante ao sol da janela para térmicas;
anemômetro de celular/ventoinha em alturas diferentes no pátio para o
gradiente do perfil (qualitativo, mas memorável).

\subsection*{Problemas sugeridos do livro (exercícios do Cap.~18)}

Aquecimento: $z_i$ por encroachment com números dados; $u(z)$ do log.
Aprofundamento: TKE com fluxos dados; jato noturno com $f$ de
latitudes diferentes. Síntese: ``por que planadores decolam às 11\,h
e pousam às 17\,h?''. \emph{Confira a numeração na sua cópia
(v1.02b).}

\subsection*{Uso do notebook \texttt{cap18\_camadalimite.ipynb}}

\begin{center}
\begin{tabular}{lll}
\hline
Caso & Conteúdo & Momento sugerido \\
\hline
1 & crescimento de $z_i$ (EDO vs.\ raiz) & Aula 1 \\
2 & ajuste do perfil log $+$ potência eólica & Aula 1 \\
3 & espectro $-5/3$ de sinal sintético (FFT) & Aula 2 \\
4 & jato de Blackadar (hodógrafo circular) & Aula 2 \\
\hline
\end{tabular}
\end{center}

O Caso~3 rende discussão de método experimental (taxa de amostragem,
aliasing — N3): meteorologia como física de laboratório.

\subsection*{Exercícios complementares e soluções}

Nas notas: T1--T5 e N1--N4. Soluções em \texttt{cap18\_solucoes.ipynb}
(\emph{não distribuir}): T1 e T4 com \texttt{sympy}; N4 explica por
que o LLJ é fenômeno subtropical/de latitudes médias. Pares:
T1$\leftrightarrow$N1, T2$\leftrightarrow$N2, T4$\leftrightarrow$N4,
T5$\leftrightarrow$N3.

\subsection*{Conexões com o resto do curso}

Balanço de superfície (Cap.~3); perfis e estabilidade (Cap.~5);
entranhamento (Cap.~14); Ekman e spin-down (Caps.~10, 13); brisas e
catabáticos como CLAs especiais (Cap.~17); dispersão de poluentes usa
TUDO deste capítulo (Cap.~19); parametrização de CLA nos modelos
(Cap.~20).
